\documentclass[11pt,a4paper]{article}
\usepackage[margin=2.5cm]{geometry}
\usepackage{amsmath,amssymb,mathtools,bm}
\usepackage{booktabs}
\usepackage{enumitem}
\usepackage[dvipsnames]{xcolor}
\usepackage[colorlinks=true,linkcolor=blue!50!black,citecolor=blue!50!black]{hyperref}

\newcommand{\bH}{\mathbf{H}}
\newcommand{\bL}{\mathbf{L}}
\newcommand{\bg}{\mathbf{g}}
\newcommand{\bx}{\mathbf{x}}
\newcommand{\bz}{\mathbf{z}}
\newcommand{\bb}{\mathbf{b}}
\newcommand{\bu}{\mathbf{u}}
\newcommand{\be}{\mathbf{e}}
\newcommand{\bd}{\boldsymbol{\Delta}}
\newcommand{\bk}{\boldsymbol{\kappa}}
\newcommand{\bG}{\mathbf{G}}
\newcommand{\bX}{\mathbf{X}}
\newcommand{\bP}{\mathbf{P}}
\newcommand{\by}{\mathbf{y}}
\newcommand{\blam}{\boldsymbol{\lambda}}
\newcommand{\bone}{\mathbf{1}}
\DeclareMathOperator{\diag}{diag}

\title{The Norm-Extended Optimization (NEO) Algorithm\\
for Second-Order Restricted-Step Orbital Optimization\\
{\large Method-independent summary: level shifting, trust region, saddle-point targeting}}
\author{Compiled from Jensen \& J\o rgensen, J.\ Chem.\ Phys.\ \textbf{80}, 1204 (1984) [JJ84];\\
Jensen, Dyall, Saue \& F\ae gri, J.\ Chem.\ Phys.\ \textbf{104}, 4083 (1996) [JDSF96];\\
Thyssen, Fleig \& Jensen, J.\ Chem.\ Phys.\ \textbf{129}, 034109 (2008) [TFJ08];\\
Nottoli, Gauss \& Lipparini, J.\ Chem.\ Theory Comput.\ (2021) [NGL21]}
\date{}

\begin{document}
\maketitle
\tableofcontents

%======================================================================
\section{Scope}
%======================================================================
This note summarizes the norm-extended optimization (NEO) algorithm as a \emph{generic} second-order, restricted-step (trust-region)
optimizer for orbital rotations. Nothing below assumes a particular electronic-structure method and nothing assumes that the energy is
invariant under any subset of rotations: the parameter vector is simply whatever set of rotation parameters the method treats as variables, all on the same footing.
Only three ingredients from the method are required, all evaluated at the current expansion point (CEP):
\begin{enumerate}[nosep]
 \item the energy $E$,
 \item the gradient $\bg$ with respect to the rotation parameters,
 \item Hessian--vector products $\bH\bb$ (analytic, or by directional differentiation of the gradient).
\end{enumerate}
The original derivation [JJ84] uses configuration amplitudes as the ``norm'' variable that is carried alongside the orbital rotations. Here that role is played by
a single auxiliary scalar; this is exactly the ``augmented Hessian'' special case discussed in [JJ84, Sec.~II.H] and it retains all NEO properties:
\begin{enumerate}[nosep]
 \item \textbf{Quadratic local convergence}, as for Newton--Raphson (NR), because the level shift that appears is $\mathcal O(\|\bg\|^2)$.
 \item \textbf{Control of the Hessian index}: the $i$-th eigenvector of the extended matrix gives a step towards a stationary point whose Hessian has exactly $i-1$ negative eigenvalues. Minima ($i=1$) and \emph{saddle points of any prescribed order} are handled by the same algorithm.
 \item \textbf{Built-in trust-region compatibility}: an extra scalar $\alpha$ rescales the level shift monotonically and allows any step length to be enforced.
 \item \textbf{Matrix-free}: only $\bH\bb$ products are needed.
\end{enumerate}

%======================================================================
\section{Setup and notation}
%======================================================================
At macro-iteration $k$ the current orbitals $C^{(k)}$ (the CEP) are updated through the unitary exponential map
\begin{equation}
 C=C^{(k)}\,e^{R(\bk)},\qquad R(\bk)^\dagger=-R(\bk),\qquad R\ \text{linear in the real parameter vector }\bk\in\mathbb R^n .
\end{equation}
Complex rotation parameters are represented by their real and imaginary parts (see Sec.~\ref{sec:complex} for the doubled complex form).
The energy as a function of the parameters is $f(\bk)=E[C^{(k)}e^{R(\bk)}]$; the CEP is $\bk=\mathbf 0$ (the orbitals are re-referenced after each step, so the expansion is always about $\mathbf 0$).
Define
\begin{equation}
 f_0=f(\mathbf 0),\qquad \bg=\nabla f(\mathbf 0)\in\mathbb R^n,\qquad \bH=\nabla\nabla^Tf(\mathbf 0)\in\mathbb R^{n\times n}.
\end{equation}
\textbf{Important:} $\bH$ is the Hessian of $f(\bk)$ in \emph{this} parametrization, including the contribution from the second-order term $R^2/2$ of the exponential.
If it is obtained as a directional derivative of the gradient evaluated at $C^{(k)}e^{R(\eta\bb)}$ this is automatically the case.
The second-order model and step are
\begin{equation}
 q(\bx)=f_0+\bg^T\bx+\tfrac12\bx^T\bH\bx,\qquad \bd=\bx\ \text{(the step in parameter space)}.
\end{equation}
Eigenvalues of $\bH$ are $\epsilon_1\le\dots\le\epsilon_n$. Nothing is assumed about them: $\bH$ may be indefinite (and this is the situation of interest for saddle points), and it may have zero eigenvalues.

%======================================================================
\section{The NEO eigenproblem}
%======================================================================

\subsection{Extended matrix}
Introduce the extended vector $\bz=(z_0;\bx)\in\mathbb R^{n+1}$ and the extended (``augmented Hessian'') matrix
\begin{equation}
 \bL(\alpha)=\begin{pmatrix}0&\alpha\,\bg^T\\ \alpha\,\bg&\bH\end{pmatrix},\qquad
 \bL\equiv\bL(1).\label{eq:L}
\end{equation}
Note that the second-order model is a Rayleigh-type form of $\bL$ on the slice $z_0=1$:
\begin{equation}
 q(\bx)-f_0=\tfrac12\,\bz^T\bL\,\bz,\qquad \bz=(1;\bx).
\end{equation}
Properties (verified directly from \eqref{eq:L}, with $\be_0=(1;\mathbf 0)$ the ``CEP vector'' and $\bG_e=(0;\bg)$):
\begin{equation}
 \bL\be_0=\bG_e,\qquad \be_0^T\bL\be_0=0 .\label{eq:Lprops}
\end{equation}
Hence, as long as $\bg\neq\mathbf 0$, $\bL$ has both positive and negative eigenvalues (at least one negative). $\be_0$ is an eigenvector of $\bL$ iff $\bg=\mathbf 0$.
In the NEO formulation of [JJ84] the corresponding statements are $\bX^T\bL\bX=0$ and $\bL\bX=\bG$, with $\bX$ the normalized CI-amplitude vector.

\subsection{Eigenproblem and the level-shifted Newton equation}
Solve for the $i$-th eigenpair (ascending order) of $\bL(\alpha)$ with $z_0\neq0$ and normalize $z_0=1$, $\bz=(1;\bx)$:
\begin{equation}
 \bL(\alpha)\bz=\lambda\bz\quad\Longleftrightarrow\quad
 \begin{cases}\alpha\,\bg^T\bx=\lambda,\\[2pt] \alpha\,\bg+\bH\bx=\lambda\bx.\end{cases}\label{eq:eigsys}
\end{equation}
The second row gives $(\bH-\lambda\bone)\bx=-\alpha\bg$, so that with the step $\bd=\bx/\alpha$
\begin{equation}
 \boxed{\ \bd=-(\bH-\lambda\bone)^{-1}\bg,\qquad \lambda=\alpha\,\bg^T\bx=\alpha^2\bg^T\bd=-\alpha^2\,\bg^T(\bH-\lambda\bone)^{-1}\bg\ }\label{eq:neostep}
\end{equation}
i.e.\ \textbf{a Newton--Raphson step with the Hessian level-shifted by the self-consistently determined shift $\lambda$}.
The last relation is a secular equation for $\lambda$ (it can be solved directly if the spectrum of $\bH$ is available, but in a matrix-free setting one just diagonalizes $\bL$ iteratively).
With $\alpha=1$ the shift satisfies $\lambda=-\bg^T(\bH-\lambda\bone)^{-1}\bg=\mathcal O(\|\bg\|^2)$; for $\lambda=0$ one recovers the NR step $-\bH^{-1}\bg$. In normalized-eigenvector language (unit-norm $\bz$) the step is $\bd=(z_0\alpha)^{-1}\bz_{1:n}$; equivalently, rescale to $z_0=1$.

\subsection{Interlacing and the Hessian index}
$\bL(\alpha)$ is $\bH$ bordered by a row and column, so by Cauchy interlacing, for every $\alpha$ (including $\alpha\to\pm\infty$),
\begin{equation}
 \lambda_1\le\epsilon_1\le\lambda_2\le\epsilon_2\le\cdots\le\epsilon_n\le\lambda_{n+1}.\label{eq:interlace}
\end{equation}
Hence $\epsilon_{i-1}\le\lambda_i\le\epsilon_i$, and the shifted matrix $\bH-\lambda_i\bone$ has \emph{exactly $i-1$ negative eigenvalues}. The NEO step is therefore a stationary point of the level-shifted quadratic model with prescribed signature: the energy is maximized along $i-1$ Hessian directions and minimized along the remaining ones. Consequently
\begin{center}
\begin{tabular}{ll}
\toprule
Target & Eigenvector of $\bL$ used\\
\midrule
Minimum & lowest ($i=1$)\\
Saddle point of order $m$ (Hessian with $m$ negative eigenvalues) & $i=m+1$\\
\bottomrule
\end{tabular}
\end{center}
The index is followed through the position of the selected eigenvalue in the spectrum of $\bL$, without diagonalizing $\bH$.
The algorithm is \emph{not} restricted to minima; only the eigenvector index and the trust-radius rule (Sec.~\ref{sec:ratio}) change for saddle points.
(A typical example is a problem where the energy must be maximized with respect to $m$ known directions and minimized with respect to all others; the target index is then $m+1$ and the converged Hessian has $m$ negative eigenvalues.)

\smallskip\noindent\emph{Caveat.} If $\bg$ is exactly orthogonal to an eigenvector of $\bH$ (e.g.\ by symmetry), the corresponding $\epsilon_j$ is also an eigenvalue of $\bL$ with $z_0=0$; such decoupled eigenvectors do not carry a step and must be deflated or counted explicitly when the index $i$ is chosen. Only eigenvectors with $z_0\neq0$ define steps.

\subsection{Quadratic convergence}
The step \eqref{eq:neostep} equals the NR step for $\bH-\lambda\bone$. Expanding $\bg$ and $\bH$ around the stationary point with error $\bm\varepsilon^{(k)}$: the NR error terms are quadratic in $\bm\varepsilon^{(k)}$, and the shift contributes terms that are of third and higher order because $\lambda=\mathcal O(\|\bg\|^2)=\mathcal O(\|\bm\varepsilon\|^2)$ (for $\alpha=1$). So the local convergence rate is quadratic, as for NR without any level shift.

\subsection{Predicted energy change}
For any $\alpha$, with $\bx$ from \eqref{eq:eigsys} ($z_0=1$) and step $\bd=\bx/\alpha$,
\begin{equation}
 \Delta q=\bg^T\bd+\tfrac12\bd^T\bH\bd=\frac{\lambda}{2\alpha^2}\big(1+\|\bx\|^2\big)=\frac{\lambda}{2}\big(\alpha^{-2}+\|\bd\|^2\big),\qquad
 \|\bd\|=\|\bx\|/\alpha .\label{eq:dq}
\end{equation}
For $\alpha=1$ and a unit-norm eigenvector $\bz$: $\Delta q=\lambda/(2z_0^2)$. (Derivation in App.~\ref{app:dq}.) This is the ``predicted'' change used in the trust-region ratio.
In practice $\Delta q$ can also be evaluated as $\bg^T\bd+\tfrac12\bd^T\bH\bd$ directly, reusing the stored $\bH\bb$ products of the Davidson subspace (since $\bd$ lies in that subspace).

%======================================================================
\section{Restricted step and level shifting}\label{sec:trust}
%======================================================================

\subsection{Trust-region subproblem}
The model is only trustworthy for $\|\bd\|\le h_k$ (trust radius). If no acceptable stationary point exists inside, the optimum lies on the boundary $\|\bd\|=h_k$, found from the Lagrangian
\begin{equation}
 \mathcal L(\bd)=f_0+\bg^T\bd+\tfrac12\bd^T\bH\bd-\tfrac{\nu}{2}(\bd^T\bd-h_k^2)\ \Longrightarrow\ 
 \bd(\nu)=-(\bH-\nu\bone)^{-1}\bg,
\end{equation}
\begin{equation}
 \Delta^2(\nu)=\bd(\nu)^T\bd(\nu)=\sum_{j=1}^n\frac{\gamma_j^2}{(\epsilon_j-\nu)^2}=h_k^2,\qquad \gamma_j=\bu_j^T\bg,\ \ \bH\bu_j=\epsilon_j\bu_j. \label{eq:steplen}
\end{equation}
$\Delta^2(\nu)$ has asymptotes at the $\epsilon_j$; the roots of $\Delta^2(\nu)=h_k^2$ are the candidate shifts. This is the Levenberg--Marquardt / restricted-step framework; NEO provides the shift through an eigenproblem, and the parameter $\alpha$ lets it be tuned.

\paragraph{Conditions on $\nu$ for a stationary point of order $i-1$.} With $\epsilon_0=-\infty$, $\epsilon_{n+1}=+\infty$:
\begin{align}
 &\text{(1)}\quad \epsilon_{i-1}<\nu<\epsilon_i,\label{eq:c1}\\
 &\text{(2)}\quad \tfrac12\epsilon_{i-1}<\nu<\tfrac12\epsilon_i.\label{eq:c2}
\end{align}
Condition (1) makes the first-order energy $E^{[1]}=f_0+\bg^T\bd$ give a positive contribution from Hessian modes $1,\dots,i-1$ (maximized) and a negative one from modes $i,\dots,n$ (minimized). Condition (2) enforces the same sign pattern for the second-order energy $q$.
For a minimum, $\nu$ is taken from the left-most branch of $\Delta^2(\nu)$ ($\nu<\epsilon_1$), which is unique and guarantees convergence. For saddles both conditions may not be simultaneously satisfiable (when $\tfrac12\epsilon_{i-1}>\epsilon_i$ or $\epsilon_{i-1}>\tfrac12\epsilon_i$); the remedy in the literature is a coordinate scaling of the variables (Simons et al.) requiring the lowest $i$ eigenvalues of $\bH$. For a minimum no scaling is needed.
For saddle points there may be no boundary stationary point with the correct structure; then a $\nu$ in the desired range is chosen and the step scaled down to the boundary.
Note that changing $\nu$ changes the direction as well as the length of the step (for shorter steps the gradient term dominates the Hessian term).

\subsection{Level shifting inside NEO: the role of $\alpha$}
With $\alpha=1$ the eigenproblem produces the shift $\lambda_i$; property \eqref{eq:interlace} guarantees condition (1) for any target index. If $\|\bd\|\le h_k$ the step is accepted (and is optimal because of quadratic convergence). Otherwise the shift is changed by rescaling the coupling to the gradient through $\alpha>1$: from \eqref{eq:neostep},
\begin{equation}
 (\bH-\lambda_i^\alpha\bone)\bd=-\bg,\qquad \lambda^\alpha_i=-\alpha^2\bg^T(\bH-\lambda_i^\alpha\bone)^{-1}\bg .
\end{equation}
Properties (Appendix~A of [JJ84], specialized here):
\begin{itemize}[nosep]
 \item Interlacing \eqref{eq:interlace} holds for every $\alpha$, so $\lambda^\alpha_i$ always lies in $[\epsilon_{i-1},\epsilon_i]$: it is the only shift in the correct domain for order $i-1$.
 \item $\alpha=0$: $\bL(0)=\diag(0,\bH)$, eigenvalues $\lambda^0=\{0\}\cup\{\epsilon_j\}$.  $\alpha=1$: plain NEO.
 \item $\alpha\to\infty$: $\bL(\alpha)\approx\alpha\begin{pmatrix}0&\bg^T\\ \bg&0\end{pmatrix}$ with eigenvalues $\pm\alpha\|\bg\|$ (eigenvectors $\propto(1;\pm\bg/\|\bg\|)$). Hence $\lambda_1^\alpha\to-\infty$, $\lambda^\alpha_{n+1}\to+\infty$, with steps along $-\bg$ (steepest descent) and $+\bg$ (steepest ascent), as expected. Intermediate roots stay bracketed by $\epsilon$'s.
 \item Monotonicity: differentiating the secular equation,
 \begin{equation}
  \alpha\frac{\partial\lambda_i^\alpha}{\partial\alpha}=2\Big[1+\alpha^2\bg^T(\bH-\lambda_i^\alpha\bone)^{-2}\bg\Big]^{-1}\lambda^\alpha_i,
 \end{equation}
 so $\lambda_i^\alpha$ is a monotonic function of $\alpha$ (its sign is that of $\lambda^\alpha_i$, and the bracket is positive) and every admissible shift corresponds to a unique $\alpha$. For the lowest state $\lambda_1^\alpha$ decreases monotonically from $\lambda_1$ to $-\infty$ as $\alpha:1\to\infty$.
\end{itemize}
\textbf{Choice of $\alpha$} ($\alpha\ge1$ suffices, since $\alpha=1$ is the quadratically convergent choice):
\begin{enumerate}[nosep]
 \item Solve with $\alpha=1$. If $\|\bd\|\le h_k$, accept.
 \item Otherwise find $\alpha>1$ such that $\|\bd_\alpha\|=\|\bx_\alpha\|/\alpha\approx h_k$ (within about $10\%$). A one-dimensional root search on $\alpha$ is enough since the map is monotonic.
\end{enumerate}
For a saddle target ($i>1$) the required $\alpha$ is the one that brings $\lambda^\alpha_i$ into the range dictated by conditions (1)--(2).
Efficient implementation: once the Davidson subspace has converged for the projected matrix, changing $\alpha$ only changes the reduced matrix (see Sec.~\ref{sec:direct}); the new eigenpairs converge in very few additional micro-iterations. The same scaling is used in [NGL21]: $\bL(\alpha)$ with step $\bx_0+\alpha^{-1}\bP\by(\alpha)$ and $\|\bx(\alpha)-\bx_0\|=R_t$.

%======================================================================
\section{Trust radius update: convergence ratio and Fletcher's algorithm}\label{sec:ratio}
%======================================================================
After a trial step, compare the true energy change with the change predicted by the second-order model. The \textbf{ratio}
\begin{equation}
 \boxed{\ r_k=\frac{E(\bx_k)-E(\bx_{k-1})}{q(\bx_k)-E(\bx_{k-1})}=\frac{\Delta E_k}{\Delta q_k}\ }
\end{equation}
with, for a step $\bd$ from the CEP,
\begin{equation}
 \Delta q=\bg^T\bd+\tfrac12\bd^T\bH\bd,\qquad \Delta E=E(\text{after step})-E(\text{CEP})=\Delta q+R^{(2)},
\end{equation}
($R^{(2)}$ the remainder beyond second order; $\Delta q$ is given by \eqref{eq:dq} for a NEO step). The closer $r_k$ is to $1$, the more quadratic the energy surface. Suggested initial radius $h_0=0.5$ ($\approx\pi/6$ for orbital-rotation angles); never accept $h_k>0.75$ ($\approx\pi/4$), since a second-order expansion cannot describe rotations larger than that.

\subsection{Minimization}
\begin{enumerate}[nosep,label=(\roman*)]
 \item $r_k<0.25$: $h_k=\tfrac23h_{k-1}$;
 \item $0.25<r_k<0.75$: $h_k=h_{k-1}$;
 \item $r_k>0.75$: $h_k=\min(1.2\,h_{k-1},\,0.75)$;
 \item $r_k<0$: \textbf{reject} the step (energy increased), and recompute a step inside a smaller radius $\tfrac23h_{k-1}$.
\end{enumerate}

\subsection{Saddle points}
Here the energy change is a cancellation of positive (maximized modes) and negative (minimized modes) contributions, so $r_k\gg1$ cannot be trusted; the acceptance window is made symmetric about $r=1$. With parameters $a_h=1.2$, $r_{\min}=0.7$, $r_{\rm good}=0.85$ (values from the literature; the algorithm is insensitive to them):
\begin{enumerate}[nosep,label=(\roman*)]
 \item $r_{\min}<r_k<r_{\rm good}$ or $2-r_{\rm good}<r_k<2-r_{\min}$: $h_k=h_{k-1}$;
 \item $r_{\rm good}<r_k<2-r_{\rm good}$: $h_k=\min(a_h\,h_{k-1},\,0.75)$;
 \item $r_k<r_{\min}$ or $r_k>2-r_{\min}$: reject the step and recompute with $h_k=h_{k-1}/a_h$.
\end{enumerate}
A rejected step is redone with the same gradient and Hessian information (only $\alpha$ changes), so back-stepping is cheap: the Davidson subspace and stored products can be reused [TFJ08].

%======================================================================
\section{Direct (matrix-free) implementation}\label{sec:direct}
%======================================================================
\subsection{Extended matrix--vector product}
For a trial vector $\bb=(b_0;\bb_x)$,
\begin{equation}
 \bL(\alpha)\bb=\begin{pmatrix}\alpha\,\bg^T\bb_x\\ \alpha\,b_0\,\bg+\bH\bb_x\end{pmatrix},
\end{equation}
so a single Hessian--vector product $\bH\bb_x$ per trial vector is needed; the rest is a scalar product and a vector addition. Changing $\alpha$ costs only $\bL(\alpha)\bb=\bL(1)\bb+(\alpha-1)\begin{pmatrix}\bg^T\bb_x\\ b_0\bg\end{pmatrix}$, i.e.\ it needs no new Hessian products. Store $\bH\bb_j$ for the subspace, then all reduced matrices for any $\alpha$ are obtained by projection:
\begin{equation}
 \bL^{\rm red}_{jl}(\alpha)=\bb_j^T\bL(\alpha)\bb_l .
\end{equation}

\subsection{Hessian--vector products}
Options, all avoiding explicit Hessian formation:
\begin{itemize}[nosep]
 \item analytic (method-dependent) one-index-transformed expressions;
 \item directional finite differences of the gradient with the \emph{same} exponential parametrization,
 \begin{equation}
  \bH\bb\simeq\frac{\bg(\eta\bb)-\bg(-\eta\bb)}{2\eta},\qquad \bg(\bk)=\nabla f(\bk)\ \text{evaluated at}\ C^{(k)}e^{R(\bk)},
 \end{equation}
 (central differences are preferable to forward differences for accuracy of the eigenvalue/index; $\eta$ is chosen small but well above round-off; all other variables of the method are held fixed).
\end{itemize}
Because the gradient is taken with respect to $\bk$ in the exponential map, the finite-difference Hessian automatically includes the second-order exponential term.

\subsection{Davidson iteration for the extended eigenproblem}
\begin{enumerate}[nosep]
 \item Trial space initialization: $\be_0=(1;\mathbf 0)$ and $(0;\bg)$ (the two starting vectors recommended in [JJ84]). For saddle/excited targets, add previously converged eigenvectors from the last macro-iteration.
 \item Reduced matrix $\bL^{\rm red}(\alpha)$, diagonalize, select the $i$-th ascending eigenpair (Ritz value $\theta_i$, coefficient vector). A block (simultaneous-expansion, Liu) Davidson with all $i$ lowest roots is recommended when $i>1$, since $\bL$ may have many small eigenvalues and single-root iterations may converge slowly. Track the root by overlap with the previous macro-iteration if needed.
 \item Residual $\mathbf r=\bL\bz-\theta_i\bz$; stop when $\|\mathbf r\|$ is below the micro tolerance; otherwise precondition (e.g.\ with the diagonal of $\bH$: $r_p/(\theta_i-H_{pp})$-type update) and add the new vector.
 \item The micro-tolerance only needs to match the accuracy of the macro-iteration (typically tied to $\|\bg\|$); 3--10 Hessian-product-evaluations per macro-iteration are common in the original applications.
\end{enumerate}

%======================================================================
\section{Complex parameters and mixed-type trial vectors}\label{sec:complex}
%======================================================================
\paragraph{Complex rotation parameters.} If the method has complex rotation parameters $\kappa_j$ (anti-Hermitian generator $R$), there are two equivalent treatments.
(a) Split into real and imaginary parts and use the real algebra above; the index counts and the trust radius refer to this real parameter space. (b) Use the doubled vector $\blam=(\bk;\bk^*)$ with gradient $\mathbf E^{[1]}=(\partial E/\partial\bk^*;\partial E/\partial\bk)$ and Hermitian Hessian
\begin{equation}
 \mathbf E^{[2]}=\begin{pmatrix}\partial^2E/\partial\bk^*\partial\bk&\partial^2E/\partial\bk^*\partial\bk^*\\ \partial^2E/\partial\bk\partial\bk&\partial^2E/\partial\bk\partial\bk^*\end{pmatrix},\qquad
 \bL(\alpha)=\begin{pmatrix}0&\alpha\,\mathbf E^{[1]\dagger}\\ \alpha\,\mathbf E^{[1]}&\mathbf E^{[2]}\end{pmatrix}
\end{equation}
(Hermitian). All statements of Secs.~3--5 hold with transposes replaced by adjoints; the doubled Hessian has the same spectrum as the real Hessian in the (real, imaginary) representation, so index counts coincide, provided the eigenvectors respect the conjugation structure $\blam=(\bk;\bk^*)$ (as in [JDSF96], where this doubled form was used).

\paragraph{Known maximization directions (pure-type trial vectors).} In problems where a known set of directions must be \emph{maximized} while all others are minimized (a saddle point of known order $m$ with identified negative-curvature subspace), it is helpful to control the minimax structure by keeping each trial vector \emph{purely} inside that subspace or purely inside its orthogonal complement [JDSF96, TFJ08]. With $p$ trial vectors of the first type, the restricted-step shift for the $s$-th target state must fall in $\epsilon^{\rm red}_{p+s-1}<\nu<\epsilon^{\rm red}_{p+s}$ (reduced Hessian eigenvalues); the NEO analogue is to select the $(p+s)$-th eigenvalue of the reduced extended matrix. This is the same as choosing $i=m+1$ once the subspace contains all $m$ maximization directions. If the negative-curvature subspace is not known in advance, use $i=m+1$ without this device.

%======================================================================
\section{Consolidated algorithm}
%======================================================================
\begin{enumerate}[label=\textbf{Step \arabic*.},leftmargin=*,nosep]
 \item \textbf{Inputs}: orbitals $C^{(k)}$; target order $m$ (number of negative Hessian eigenvalues at the solution; $m=0$ for minima), eigenvector index $i=m+1$; trust radius $h_k$; thresholds.
 \item Evaluate $E$ and $\bg$ at the CEP. If $\|\bg\|<$ tolerance: converged; verify that the final Hessian (or the position of the converged root) corresponds to order $m$.
 \item If $k>1$: compute $r_{k-1}=\Delta E/\Delta q$ and update $h_k$ (or reject the last step and back-step) using Sec.~\ref{sec:ratio}.
 \item Micro-iterations: Davidson on $\bL(\alpha=1)$ (Sec.~\ref{sec:direct}), selecting the $i$-th ascending eigenvalue; set $\bx$ from $z_0=1$, step $\bd=\bx$.
 \item If $\|\bd\|>h_k$: increase $\alpha$ (monotonic root search on the reduced matrix) until $\|\bd_\alpha\|=\|\bx_\alpha\|/\alpha\approx h_k$; $\bd=\bx_\alpha/\alpha$.
 \item Compute $\Delta q=\bg^T\bd+\tfrac12\bd^T\bH\bd$ from stored products.
 \item Update orbitals: $C^{(k+1)}=C^{(k)}\,e^{R(\bd)}$; set $k\leftarrow k+1$; go to Step 2.
\end{enumerate}

%======================================================================
\section{Practical checks}
%======================================================================
\begin{itemize}[nosep]
 \item Model consistency: $\be_0^T\bL\be_0=0$ and $\bL\be_0=(0;\bg)$; eigenvalue interlacing \eqref{eq:interlace} against a dense $\bH$ on small test cases.
 \item Step identity: the NEO step equals $-(\bH-\lambda_i\bone)^{-1}\bg$ with $\lambda_i=\alpha\bg^T\bx$; and $\bH-\lambda_i\bone$ has $i-1$ negative eigenvalues.
 \item Monotonic behaviour of $\lambda_i^\alpha$ and $\|\bd_\alpha\|$ with $\alpha$.
 \item Predicted change \eqref{eq:dq} equals $\bg^T\bd+\tfrac12\bd^T\bH\bd$.
 \item Gradient/Hessian consistency: finite difference of the energy versus $\bg$, and of $\bg$ versus $\bH\bb$, using the exponential map.
 \item Quadratic decrease of $\|\bg\|$ for $\alpha=1$ near convergence; if it is only linear, the Hessian--vector products (or the exponential-map second-order term) are inconsistent with the gradient.
 \item Relation to trust-region Newton--CG (Steihaug--Toint): both solve the same trust-region model; NEO obtains the boundary step and the index information from an eigenproblem, while truncated CG only detects negative curvature and terminates on the boundary; for prescribed-index saddle points, NEO gives direct control through $i$ and $\alpha$.
\end{itemize}

%======================================================================
\appendix
\section{Derivations}
%======================================================================
\subsection{Predicted energy change}\label{app:dq}
From \eqref{eq:eigsys} with $z_0=1$: $\alpha\bg^T\bx=\lambda$ and $\bH\bx=\lambda\bx-\alpha\bg$. Then
\begin{equation}
 \bx^T\bH\bx=\lambda\|\bx\|^2-\alpha\bg^T\bx=\lambda\|\bx\|^2-\lambda .
\end{equation}
With $\bd=\bx/\alpha$:
\begin{equation}
 \Delta q=\bg^T\bd+\tfrac12\bd^T\bH\bd=\frac{\lambda}{\alpha^2}+\frac{\lambda\|\bx\|^2-\lambda}{2\alpha^2}=\frac{\lambda}{2\alpha^2}\big(1+\|\bx\|^2\big).
\end{equation}
For a unit-norm eigenvector $\bz$, $\|\bz\|^2=z_0^2(1+\|\bx\|^2)=1$, so at $\alpha=1$: $\Delta q=\lambda/(2z_0^2)$.

\subsection{Secular equation and monotonicity in $\alpha$}
From $(\bH-\lambda\bone)\bd=-\bg$ and $\lambda=\alpha^2\bg^T\bd$: $\lambda=-\alpha^2\bg^T(\bH-\lambda\bone)^{-1}\bg$. Differentiating with respect to $\alpha$, using $\partial_\alpha(\bH-\lambda\bone)^{-1}=(\bH-\lambda\bone)^{-2}\,\partial_\alpha\lambda$:
\begin{equation}
 \partial_\alpha\lambda=\frac{2\lambda}{\alpha}-\alpha^2\bg^T(\bH-\lambda\bone)^{-2}\bg\ \partial_\alpha\lambda
 \ \Longrightarrow\
 \alpha\,\partial_\alpha\lambda=\frac{2\lambda}{1+\alpha^2\bg^T(\bH-\lambda\bone)^{-2}\bg}.
\end{equation}
The denominator is positive, so $\alpha\partial_\alpha\lambda$ has the sign of $\lambda$; for the lowest root ($\lambda<0$) $\lambda$ decreases monotonically with $\alpha>0$.

\subsection{Step-length function}
Expanding $\bg=\sum_j\gamma_j\bu_j$ in the eigenvectors of $\bH$: $\bd(\nu)=-\sum_j\gamma_j(\epsilon_j-\nu)^{-1}\bu_j$, hence $\|\bd(\nu)\|^2=\sum_j\gamma_j^2/(\epsilon_j-\nu)^2$, which is increasing in $|\nu-\epsilon_j|^{-1}$ near each asymptote and monotonic on the branch $\nu<\epsilon_1$; this branch has a single root of $\|\bd\|=h_k$ (for $\bg$ not orthogonal to the lowest eigenvector), and corresponds to the minimum-seeking trust-region step.

\end{document}
